% V17 independent working edition. V16 retains the time-driven model.
\IfFileExists{siamonline250211.cls}{%
 \documentclass[review,hidelinks,onefignum,onetabnum]{siamonline250211}%
 \def\isSIAM{1}%
}{%
 \documentclass[11pt]{article}%
 \def\isSIAM{0}%
}
% V17 uses the V16 SIAM manuscript interface. The article fallback is solely
% for offline preview when the repository's unchanged SIAM ZIP is unavailable.
\usepackage{microtype,xurl,needspace}
\usepackage{amsmath,amsfonts,amssymb,amsthm,mathtools,bm}
\usepackage{graphicx,booktabs,enumitem,tabularx}
\usepackage[numbers,sort&compress]{natbib}
\ifnum\isSIAM=1
  \newsiamremark{assumption}{Assumption}
  \newsiamremark{remark}{Remark}
  \headers{When Can Markets Crash Without News?}{Siyuan Ye}
  \renewcommand{\footercopyright}{}
  \setlength{\headheight}{14pt}
\else
  \usepackage[hidelinks]{hyperref}
  \usepackage[margin=1in]{geometry}
  \newtheorem{theorem}{Theorem}[section]
  \newtheorem{lemma}[theorem]{Lemma}
  \newtheorem{proposition}[theorem]{Proposition}
  \newtheorem{corollary}[theorem]{Corollary}
  \newtheorem{definition}[theorem]{Definition}
  \newtheorem{assumption}[theorem]{Assumption}
  \newtheorem{remark}[theorem]{Remark}
  \newcommand{\headers}[2]{}
  \newenvironment{keywords}{\noindent\textbf{Keywords.} }{\par}
  \newenvironment{MSCcodes}{\noindent\textbf{MSC codes.} }{\par}
\fi
\title{When Can Markets Crash Without News?\\Activity-Driven Recovery and Systemic Reachability}
\author{Siyuan Ye\thanks{Independent Researcher.}}
\hypersetup{pdftitle={When Can Markets Crash Without News? Activity-Driven Recovery and Systemic Reachability},pdfauthor={Siyuan Ye}}
\setlength{\emergencystretch}{2em}
\newcommand{\dd}{\,\mathrm{d}}
\newcommand{\Cris}{\mathcal{C}}
\newcommand{\one}{\mathbf{1}}
\newcommand{\pos}[1]{\left[#1\right]_+}
\newcommand{\TV}{\operatorname{TV}}

\begin{document}
\raggedbottom
\maketitle
\begin{abstract}
Matched-dollar transfers can produce systemic terminal price losses even when
fundamental benchmarks are fixed and every asset's cumulative executed-dollar
ledger closes at zero. We replace recovery in physical time by recovery driven
exclusively by gross executed-dollar activity. The resulting nonlinear control
system has no price motion without transactions. A nine-edge three-asset
protocol establishes local all-asset declines for arbitrary smooth positive
response functions, while scalar activity floors and execution-resource
inequalities restrict severe losses. In the constant-impact model, a
continuous-direction argument provides analytic existence constructions for
all ten prices to end at half their benchmarks, and for nine prices to become
arbitrarily small with uniform finite turnover and inventory budgets.
A separate internally interval-checked, no-instantaneous-netting ten-asset
path makes five assets close at the lower limit of successive $\pm25\%$
rolling sessions. The path has finite gross share-execution rate and an
equal-weight geometric-index loss above $16\%$. These are deterministic
reachability results, not equilibrium crash predictions.
\end{abstract}
\begin{keywords}market crashes; trade activity; price impact; systemic drawdown; reachability; interval verification; execution constraints\end{keywords}
\begin{MSCcodes}91G15, 91G45, 93B03, 65G30\end{MSCcodes}

\section{Introduction}\label{sec:intro}
Crashes may accompany public news, forced selling, liquidity withdrawal or
information revealed by trades. Their respective theoretical mechanisms are
studied, among others, by \citet{GennotteLeland1990,Romer1993},
\citet{HuangWang2009,CespaFoucault2014}, and \citet{ContWagalath2013}.
An independent mathematical question is whether a price decline is compatible
with a fully paired stream of trades whose \emph{cumulative dollar balance
in each asset} is zero. This is a reachability question, not an explanation
of any particular historical crash.

Resilient execution models typically assign time dynamics to the liquidity
state; see \citet{ObizhaevaWang2013,AlfonsiSchied2010}. We instead ask what
changes when price resilience is triggered only by gross dollar activity
associated with actual trades. Such activity can be strictly positive even
when a node's simultaneous net executed-dollar flow vanishes. It is therefore
essential to distinguish a permissive control class from a stricter class
in which each node can only buy or only sell at a given instant. Our model
is an explicitly parameterized price-response law, not a strategic equilibrium
or a no-manipulation result of the type considered by
\citet{Gatheral2010,HubermanStanzl2004}.

We provide a self-contained analytic reachability chain and one narrow
computer-assisted construction. The universal results are local triangular
reachability, scalar floors, an aggregate full-coverage barrier, and necessary
turnover and inventory bounds. The ten-asset results are: (I) all ten
terminal prices equal one half of the benchmark; (II) nine losses have no
finite severity ceiling while time, turnover and inventory are bounded by
one set of constants; and (III) the same five assets close at the lower limits
of two consecutive $\pm25\%$ sessions under the stricter no-instantaneous-
netting class. The first two are analytic existence constructions whose
parameters need not coincide with the old physical-time model's numerical
budgets. The third has a reproducible, internally outward-rounded certificate.
No earlier fixed-version manuscript or its numerical certificate is
represented as a certificate for the present dynamics.

\section{The activity-driven paired-dollar market}\label{sec:model}
Fix $N\ge2$, a fixed benchmark $F>0$, positive transaction prices $S_i$,
and a positive locally $C^2$ response $g$. The complete bidirectional graph
without self-edges has incidence matrix $B$: for an edge from seller $i$ to
buyer $j$, $B_{ie}=-1$ and $B_{je}=1$. A control is a vector of nonnegative
edge-flow measures $\dd z_e$. The \emph{executed-dollar ledger} and \emph{gross
dollar activity} of asset $i$ are
\begin{equation}\label{eq:model}
 \dd X_i=\sum_e B_{ie}\dd z_e,\qquad
 \dd A_i=\sum_e|B_{ie}|\dd z_e,\qquad
 \boxed{\dd S_i=\frac{g(S_i)}{S_i}\dd X_i-\kappa(S_i-F)\dd A_i},\quad\kappa>0.
\end{equation}
The sums $\sum_i\dd X_i=0$, $\sum_i\dd A_i=2\sum_e\dd z_e$, and the
nodewise inequality $|\dd X_i|\le\dd A_i$ follow from incidence accounting.
A direct edge transfer of dollar mass $z>0$ is integrated in its own amount
parameter $a\in[0,z]$, including \emph{both} terms in \eqref{eq:model},
with $\dd\theta_i=\dd X_i/S_i$ for share positions. Finite sequences of such
flows and absolutely continuous rates are admissible while all intermediate
prices are positive. This flow-completed convention preserves the usual
chain rule; an old rule in which recovery freezes during a trade is \emph{not}
used here. If $\dd z=0$, prices do not move.

\begin{definition}[Controls and execution resources]\label{def:controls}
In $\mathsf C_0$ every nonnegative edge flow is allowed, including simultaneous
inflow and outflow at one vertex. In $\mathsf C_1\subset\mathsf C_0$, each
asset is a pure dollar buyer or a pure dollar seller almost everywhere on
active periods; equivalently $|\dot X_i|=\dot A_i$ when $\dot A_i>0$.
Switching direction over time is permitted in both classes.
All paths start at $S_i(0)=F$ and impose $X_i(T)=0$ for \emph{every} asset.
The paired turnover and cumulative sales in shares are
\begin{equation}\label{eq:resource}
 M=\sum_e\int_0^T\dd z_e=\frac12\sum_i A_i(T),\qquad
 H_i^-=\sum_{e:\,B_{ie}=-1}\int_0^T\frac{\dd z_e}{S_i};
\end{equation}
when control is absolutely continuous the gross share-execution rate is
$r_i^{\rm gross}=\dot A_i/S_i$. In $\mathsf C_0$, $H_i^-$ counts
\emph{actual gross share sales} even when simultaneous purchases cancel
net executed dollars; the net negative variation of $X_i$ would undercount
those sales. Starting inventory $H_i\ge H_i^-$ prevents short sales.
Dollar closure does not require share-flatness.
\end{definition}

For power responses let
\begin{equation}\label{eq:power}
 g_p(S)=\eta(S/F)^p,\quad D_F=F^2/\eta,\quad
 x_i=S_i/F,\quad u_i=X_i/D_F,\quad a_i=A_i/D_F,\quad K=\kappa D_F.
\end{equation}
The uniquely relevant normalized dynamics are
\begin{equation}\label{eq:normalized}
 \boxed{\dd x_i=x_i^{p-1}\dd u_i-K(x_i-1)\dd a_i},\qquad
 |\dd u_i|\le\dd a_i,\quad\sum_i\dd u_i=0.
\end{equation}
In the power-law sections, dollars, shares, and prices are explicitly
normalized by $F=\eta=1$ where indicated. The duration $T$ does not affect
unconstrained activity paths: any absolutely continuous path can be
reparameterized over any positive interval, but gross-rate limits do depend
on physical time.

\begin{definition}[Terminal systemic crisis]\label{def:crisis}
Write $L_i(T)=\log(F/S_i(T))$, $\bar L=N^{-1}\sum_iL_i$ and
$k=\lceil\alpha N\rceil$. If $Q_\alpha$ is the $k$th largest $L_i$, define
\begin{equation}\label{eq:crisis}
 \Cris_\alpha(T)=\min\{\bar L(T),Q_\alpha(T)\}.
\end{equation}
The terminal event is $\Cris_\alpha(T)\ge\delta>0$; neither rolling
peak drawdowns nor first-passage times are identified with this criterion.
\end{definition}

\subsection{Exact transformed accounting}
Define for positive prices
\begin{equation}\label{eq:coords}
 \Phi(S)=\int_F^S\frac{v}{g(v)}\dd v,\quad
 \Psi(S)=\int_F^S\frac{\dd v}{g(v)},\quad
 E(S)=\int_F^S\frac{v-F}{g(v)}\dd v\ge0,\quad
 R(S)=\frac{S(S-F)}{g(S)}.
\end{equation}
\begin{proposition}[The activity balance and execution cost]\label{prop:account}
For each flow-completed path,
\begin{align}
 \dd\Phi(S_i)&=\dd X_i-\kappa R(S_i)\dd A_i,\label{eq:phi}\\
 \dd\left[\sum_i\Phi(S_i)\right]&=-\kappa\sum_iR(S_i)\dd A_i,\label{eq:potential}\\
 X_i(T)&=F\,[\theta_i(T)-\theta_i(0)]+E(S_i(T))
  +\kappa\int\frac{(S_i-F)^2}{g(S_i)}\dd A_i.\label{eq:cost}
\end{align}
In particular, exact dollar closure and equal share positions both at $0$
and $T$ force a stationary price path. If $X_i(T)=0$ and $S_i(T)<F$,
then $\int_{S_i>F}R(S_i)\dd A_i>0$.
\end{proposition}
\begin{proof}
Use $\Phi'(S)=S/g(S)$ in \eqref{eq:model} to get \eqref{eq:phi};
paired transfers cancel in \eqref{eq:potential}. For the last identity,
$\dd X_i=S_i\dd\theta_i$ and
$\dd X_i-F\dd\theta_i=(S_i-F)\dd X_i/S_i
=(S_i-F)\dd S_i/g(S_i)+\kappa(S_i-F)^2\dd A_i/g(S_i)$.
Integrate and apply the chain rule to $E$.
The last two nonnegative terms in \eqref{eq:cost} must both vanish under
simultaneous closures. For $S_i(T)<F$, \eqref{eq:phi} with $X_i(T)=0$
implies $-\Phi(S_i(T))=\kappa\int R(S_i)\dd A_i>0$;
$R$ is positive only above $F$.
\end{proof}

\section{Local all-asset reachability}\label{sec:local}
For $g_p$ set
\begin{equation}\label{eq:w}
 w_p(x)=\int_1^x s^{1-p}\dd s
 =\begin{cases}(x^{2-p}-1)/(2-p),&p\ne2,\\ \log x,&p=2,\end{cases}
 \qquad r_p(w_p(x))=x^{1-p}(x-1).
\end{equation}
Then $\dd w_i=\dd u_i-Kr_p(w_i)\dd a_i$ with
$r_p(0)=0$ and $r_p'(0)=1$. The same linearization works for a general
smooth positive $g$ after scaling executed dollars by $F^2/g(F)$.

\begin{theorem}[Triangular local reachability]\label{thm:triangle}
On the complete bidirectional graph, for every $N\ge3$, $\kappa>0$
and every positive $C^2$ response $g$ near $F$, a finite $\mathsf C_1$
sequence has strictly positive intermediate prices,
$X_i(T)=0$ for every $i$, and $S_i(T)<F$ for every $i$.
For a normalized power response its prices satisfy
$w_i(T)=-m_iKz^2+O_N(z^3)$ as $z\downarrow0$, with
$m_1=m_2=N-2$ and $m_j=1$ for $j\ge3$.
\end{theorem}
\begin{proof}
On a triple $(a,b,c)$ execute nine equal-mass oriented edges, in order:
\[
 a\to b,\ b\to c,\ c\to a,\ b\to c,\ c\to a,\ a\to b,\ c\to a,\ a\to b,\ b\to c.
\]
Each vertex receives and sends exactly $3z$. Its six signs in its own
activity clock are respectively
$(-,+,+,-,+,-)$, $(+,-,-,+,+,-)$, and $(+,-,+,-,-,+)$.
All three satisfy $\sum_j\sigma_j=0$ and $\sum_jj\sigma_j=-1$.
In the transformed coordinate $r(w)=w+O(w^2)$, hence variation of
constants gives
\begin{align*}
 w(6z)&=\sum_{j=1}^6\sigma_j
 \int_{(j-1)z}^{jz}e^{-K(6z-a)}\dd a+O(z^3)\\
 &=Kz^2\sum_{j=1}^6j\sigma_j+O(z^3)=-Kz^2+O(z^3).
\end{align*}
Repeat the same triple protocol on $(1,2,j)$ for $j=3,\ldots,N$.
The shared vertices start the later protocols only $O(z^2)$ away from
benchmark, so subsequent increments retain the displayed quadratic
term; $m_i$ counts participation. Each edge is a one-directional sale/purchase
at the instant of execution, hence it is in $\mathsf C_1$. For sufficiently
small $z$ the full path stays in a compact positive price neighborhood;
initial inventories covering its finite sales suffice.
\end{proof}

\begin{remark}
The distinction between $N\ge3$ and $N=2$ matters. If $N=2$ then both
assets have identical activity clocks. If $e(S)=Sg'(S)/g(S)\ge1$ on the path,
concavity of $\Phi$ gives $R(S)\le\Phi(S)$ and
$\dd H\ge-\kappa H\dd A$, where $H=\sum_i\Phi(S_i)$ and $A$ is the common
clock. Since $H(0)=0$, simultaneous losses of both assets are impossible.
If $e\ge2$, concavity of $q\mapsto\Phi(Fe^q)$ also excludes
$\Cris_\alpha>0$ by Jensen's inequality. Neither is a general
$N\ge3$ obstruction in this activity model.
\end{remark}

\section{Floors, aggregate barriers and execution resources}\label{sec:resources}
\begin{theorem}[Scalar activity-driven price floors]\label{thm:floor}
For the power response $g_p$, every individual normalized price is
bounded from below by the largest root $r\in(0,1)$ of
$K(1-r)=r^{p-1}$ whenever such a root exists. More precisely:
(i) if $p<1$, a strictly positive floor occurs for
\begin{equation}\label{eq:kc}
 K\ge K_c(p)=\frac{(2-p)^{2-p}}{(1-p)^{1-p}};
\end{equation}
(ii) for $p=1$ it occurs if $K>1$, and equals $1-1/K$;
(iii) for $p>1$ it occurs for every $K>0$.
For $p=0$, $K_c=4$ and the upper root is
$r=(1+\sqrt{1-4/K})/2$.
\end{theorem}
\begin{proof}
Using activity as clock and $|\dd u/\dd a|\le1$, the lowest possible
drift in $(0,1)$ is
$\dd x/\dd a\ge K(1-x)-x^{p-1}=:f_-(x)$.
A root of $f_-$ is a one-sided barrier by scalar comparison and local
ODE uniqueness. For $p<1$, minimize
$x^{p-1}/(1-x)$ on $(0,1)$ to get \eqref{eq:kc}, including the
noncrossable double root at equality. For $p=1$ solve the linear
root equation. For $p>1$, $f_-$ decreases strictly from $K$ at zero
to $-1$ at one and has a unique positive zero.
\end{proof}

\begin{theorem}[Necessary full-coverage bound]\label{thm:aggregate}
For any finite-variation flow-completed control in $\mathsf C_0$,
$p<2$ and $x_i(0)=1$,
\begin{equation}\label{eq:aggregate}
 \sum_{i=1}^N x_i(t)^{2-p}\ge1.
\end{equation}
In particular all $N$ prices cannot simultaneously be no larger than
$e^{-\delta}$ unless $\delta\le(\log N)/(2-p)$.
\end{theorem}
\begin{proof}
Put $\nu=2-p>0$ and $H=\sum_i w_p(x_i)$.
At every state where $H\le(1-N)/\nu$, we have
$\sum_i x_i^\nu\le1$, whence every $x_i\le1$ and
$\dd H=-K\sum_i x_i^{1-p}(x_i-1)\dd a_i\ge0$.
Initially $H=0>(1-N)/\nu$; continuity and the inward-pointing
inequality forbid crossing below $(1-N)/\nu$.
The event-flow convention permits the same argument during every edge
of a finite-variation path. Rearrangement gives \eqref{eq:aggregate}.
\end{proof}

\begin{theorem}[Necessary resources]\label{thm:lower}
Suppose at least $k$ assets end at $S_i(T)\le s=Fe^{-\delta}<F$.
Then for any $\mathsf C_0$ control,
\begin{equation}\label{eq:bounds}
 M\ge k[-\Phi(s)],\qquad
 \sum_iH_i^-\ge k[-\Psi(s)].
\end{equation}
If $X_i(T)=0$ and initial shares $H_i$ are available, then
$E(S_i(T))\le F H_i$. If $R(S)\le R_{\max}$ for all $S\in[F,U]$
and $S_i(t)\le U$ throughout, a terminally dollar-closed asset ending
below $s$ must have above-benchmark activity
\begin{equation}\label{eq:above}
 \int_{S_i>F}\dd A_i\ge\frac{-\Phi(s)}{\kappa R_{\max}}
 \quad (R_{\max}>0).
\end{equation}
\end{theorem}
\begin{proof}
On the last downward excursion from $F$ to $s$,
$\dd\Phi=\dd X-\kappa R\dd A\ge\dd X$ since $R<0$ below $F$.
Thus the negative-dollar variation of each target is at least
$-\Phi(s)$; summing sellers counts each paired dollar transfer once.
Similarly, a downward price excursion below $F$ satisfies
$\dd S\ge-g(S)(-\dd\theta)_+$ and therefore requires at least
$-\Psi(s)$ sold shares. Cost identity \eqref{eq:cost} and
$\Delta\theta_i\ge-H_i$ give the inventory inequality.
Finally, $-\Phi(S_i(T))=\kappa\int R\dd A$
under dollar closure and its positive part can only arise above $F$;
there $R\le R_{\max}$. If $S_i(T)\le s$, the stated bound follows.
\end{proof}

\begin{theorem}[A $K$-dependent turnover lower bound]\label{thm:kresource}
For $p=0$, define $D_F=F^2/\eta$. If $k$ dollar-closed assets
end at or below $s=Fe^{-\delta}$, then
\begin{equation}\label{eq:Kbound}
 M\ge kD_F\sqrt{\frac{1-e^{-2\delta}}{2K(2+K)}}.
\end{equation}
In particular this necessary lower bound diverges as $K^{-1/2}$
when $K\downarrow0$ at a fixed positive target severity.
\end{theorem}
\begin{proof}
Normalize $F=\eta=1$ and use the $i$th node's activity clock
$a\in[0,2b]$. Dollar closure makes its total buys and sells equal $b$,
so its signed ledger obeys $u(0)=u(2b)=0$ and
$u(a)\le\min(a,2b-a)$. For $w=(x^2-1)/2$ and $r=x(x-1)$,
$\dd w=\dd u-Kr\dd a$. Since $r\ge-1/4$ and
$r\le2[w]_+$ for every $x>0$,
\[
 [w(a)]_+\le\min(a,2b-a)+Ka/4,
 \qquad \int_0^{2b}r\dd a\le(2+K)b^2.
\]
At the terminal time, $-w(T)=K\int r\dd a\ge(1-s^2)/2$.
Hence $b\ge\sqrt{(1-s^2)/[2K(2+K)]}$.
The sum of all dollar purchases equals paired turnover $M$.
Restore the executed-dollar unit $D_F$.
\end{proof}

\begin{proposition}[All-node second-session obstruction]\label{prop:second}
If all assets stay at or below $F$ throughout a second trading session,
they cannot all close that session strictly below their respective
opening prices, regardless of execution-rate bounds.
\end{proposition}
\begin{proof}
The activity potential obeys \eqref{eq:potential}, and every
$R(S_i)\le0$ throughout the session, hence $H$ cannot decrease.
Strict decline of every price would strictly decrease
$H=\sum_i\Phi(S_i)$, a contradiction.
\end{proof}

\section{Three ten-asset constructions}\label{sec:examples}
We now normalize $F=\eta=1$ and use the complete graph. A useful
completion fact separates the design of a price path from the design
of a paired transaction network.
\begin{lemma}[Equal-activity graph completion]\label{lem:completion}
Let $N\ge2$, $a_0>0$, $\sum_iu_i(t)=0$ and
$|u_i(t)|\le a_0$ almost everywhere. In $\mathsf C_0$ there exist
nonnegative edge-flow rates with
$\dot X_i=u_i(t)$ and $\dot A_i=a_0$ for every $i$.
\end{lemma}
\begin{proof}
Set desired node purchases $b_i=(a_0+u_i)/2$ and sales
$s_i=(a_0-u_i)/2$. Their respective totals are $Na_0/2$.
A seller $i$ can reach every buyer except itself, and the singleton
Hall condition is $s_i+b_i=a_0\le Na_0/2$. A group of at least
two distinct sellers reaches every buyer, so every remaining
max-flow/min-cut inequality holds. A nonnegative zero-diagonal
transportation matrix therefore realizes the rates.
\end{proof}

\begin{lemma}[Continuous embedding with actual trade activity]\label{lem:embed}
Let $x_i(t)>0$ be an absolutely continuous trajectory satisfying
$\sum_i x_i(t)[\dot x_i(t)+\rho(x_i(t)-1)]=0$ and
$|u_i(t)|\le a_0$, where
$u_i=x_i[\dot x_i+\rho(x_i-1)]$ and $\rho>0$.
Then under $p=0$, $K=\rho/a_0$ and controls in $\mathsf C_0$,
$x_i$ solves the \emph{new} activity system with exactly
$\dot A_i=a_0$ and $\dot X_i=u_i$.
\end{lemma}
\begin{proof}
Apply Lemma~\ref{lem:completion}; in the new equation
$\dot x_i=u_i/x_i-K(x_i-1)a_0
=u_i/x_i-\rho(x_i-1)$, identically.
No free decay occurs: the last term is generated by the paired
gross execution flow. The turnover is $Na_0T/2$.
\end{proof}

\subsection{Construction I: all ten prices halve}\label{sec:conI}
\begin{theorem}[Ten simultaneous 50\% terminal losses]\label{thm:tenhalf}
There exists a fixed $K>0$ and a positive-price, continuous-time
$\mathsf C_0$ ten-asset path on $T=40$ with
$X_i(40)=0$ and $x_i(40)=1/2$ for all $i$. Its paired turnover
and required initial inventories are finite. This is an existence
claim, not an assertion that old-model turnover values are reproduced.
\end{theorem}
\begin{proof}
We give the explicit construction and nondegeneracy argument in
Appendix~\ref{app:smoothing}. Start from ten unit vectors $v^{(j)}$
obtained by permuting $(a,b,\ldots,b)$,
$a=\sqrt{1-9b^2}$, $0<b<1/\sqrt{10}$.
With fixed directions and instantaneous \emph{limiting} switches,
put $x_i(t)=R(t)v_i^{(j)}$, where
$R(t)=H+(\sqrt{10}-H)e^{-.4t}$, $H=a+9b$.
The equation $R(40)=\sqrt{10}/2$ has a unique interior root $b=b_0$,
because $H$ increases continuously from $1$ to $\sqrt{10}$.
The difference between a high and a low asset's accounting integrand is
\[
 W_b(t)=(a-b)R(t)[(a+b)R(t)-1]>0;
\]
indeed $(a+b)H=1+10ab>1$. Partition $[0,40]$ into ten intervals
of equal $W_b$-integral. Every asset then has the same
$\int x_i(x_i-1)\dd t$, and equal terminal prices imply equal dollar
ledgers. Aggregate dollar pairing forces each ledger to zero.

These limiting switches are \emph{not} accepted as legal new-model
impulses. Replace each by a positive-duration linear interpolation
of the squared direction components and solve the exact radial ODE
$\dot R=-.4(R-\sum_i v_i(t))$. The endpoint and nine ledger
differences define ten continuous equations in $b$ and the nine
switch times. At zero transition width their Jacobian is block
triangular with first diagonal entry
$(1-e^{-16})H'(b_0)/\sqrt{10}>0$ and cut block determinant
$10\prod_{j=1}^9 W_{b_0}(t_j)>0$.
Appendix~\ref{app:smoothing} proves the required $C^1$ convergence
of the smoothed equation map. Stability of a nondegenerate zero
thus gives an \emph{exact} positive-width solution, not a
floating-point root. For fixed positive width the net execution rates
are uniformly bounded, so Lemma~\ref{lem:embed} implements the
resulting continuous path with a fixed small $K>0$, finite turnover
and finite sales inventories. Positivity follows from the strictly
positive direction components.
\end{proof}

\subsection{Construction II: nine arbitrarily severe terminal losses}
\label{sec:conII}
\begin{theorem}[Uniform finite resources at partial coverage]\label{thm:nine}
There exist constants $T>0$, $K>0$, $M_0<\infty$ and
$H_0<\infty$, \emph{independent of} $\delta$, such that for
every finite $\delta>0$ a positive-price ten-asset $\mathsf C_0$
execution path satisfies $X_i(T)=0$ for every asset,
$M\le M_0$, $H_i^-\le H_0$, and
$\Cris_{0.9}(T)>\delta$.
\end{theorem}
\begin{proof}
Details of the uniform transition and inventory estimates are in
Appendix~\ref{app:smoothing}. Put $m=9$ and fix a small $c>0$.
Fix $c=1/100$, $T=40$, $b_-=1/10000$ and $\beta=1/5$.
During stage $j\in\{1,\ldots,9\}$ use a unit direction
$v^{(j)}=(c,\ldots,a,\ldots,c,b)$ with high entry at $j$,
$a=\sqrt{1-b^2-8c^2}$, $H=a+8c+b$ and
$b\in[b_-,\beta]$. On this entire explicit interval
$a>0.97>\max(b,c)$, $bH<1$, and $(a+c)H>1$. The limiting stage radius is
$R_b(t)=H+(\sqrt{10}-H)e^{-t}$ and the high-minus-low
integrand weight is
$W_b=(a-c)R_b[(a+c)R_b-1]>0$.
For a terminal direction
$v^{\rm fin}=(q,\ldots,q,\sqrt{1-9q^2})$, $q>0$, set
\begin{equation}\label{eq:compF}
 F_q(b)=\frac{R_b(T)^2(1-9q^2)-1}{2}
  +\int_0^T bR_b(t)[bR_b(t)-1]\dd t.
\end{equation}
The endpoint signs can be proved with explicit estimates. At
$b=b_-$, $H>1.079$ and $R(t)<4$, so
$F_0(b_-)>(1.079^2-1)/2-40(4/10000)>0$.
At $b=\beta$, $5/4<H<13/10$ and
$R(t)<41/20$ for $t\ge1$, because
$\sqrt{10}-H<2$ and $e^{-1}<3/8$. Thus
$1/4<bR(t)<41/100$ for $1\le t\le40$, while
$bR(t)<1$ for the full interval. It follows that
$F_0(\beta)<9/2-39(59/400)<0$. Hence the
compensator ledger has opposite strict signs at $b_-$ and $\beta$. Partition the nine stages into
intervals of equal $W_b$-integral, for each $b\in[b_-,\beta]$.
All nine target dollar ledgers are equal, and \eqref{eq:compF}
has an interior zero for all sufficiently small $q>0$.

Replace each instantaneous limiting switch with an interpolation
that is \emph{linear in squared directions}, including the final
transition to $v^{\rm fin}$. The eight equations making the target
accounting integrals equal have, at zero width, cut-time Jacobian
determinant $9\prod_{j=1}^8 W_b(t_j)>0$.
On the compact $b$-interval the determinants and stage gaps are
uniformly positive. A parameter-uniform inverse-function argument
therefore supplies exact cuts for all sufficiently small fixed
positive width $h$, all sufficiently small $q\ge0$, and all
$b\in[b_-,\beta]$. The compensator's ledger retains opposite
strict endpoint signs, so an intermediate $b=b(q)$ closes it.
All other ledgers close by pairing and equality of the target ledgers.

The radial bounds $1\le R\le\sqrt{10}$ and the square-direction
interpolants imply a net executed-dollar rate bounded by a constant
$a_0$ \emph{independent of $q$}; take $K=1/a_0$ and apply
Lemma~\ref{lem:embed}. This yields the same finite turnover
$M_0=5a_0T$ for every $q$. A strictly positive lower direction
$d=\min\{c,b_-,1/\sqrt{10}\}$ holds until the last transition.
During that last width-$h$ transition, each price is bounded below
by $d\sqrt{1-\lambda}$ with $\lambda\in[0,1]$.
Thus $\int_0^T\dd t/x_i(t)\le(T+h)/d$ uniformly in $q$.
Since sales dollar rate never exceeds $a_0$, all
$H_i^-\le H_0=a_0(T+h)/d$.
Finally the nine targets end at $qR(T)\le\sqrt{10}q$, while the
compensator ends between positive constants, so the mean and
ninth-largest log losses diverge as $q\downarrow0$.
Choose sufficiently small $q>0$ for each requested $\delta$.
\end{proof}

\subsection{Construction III: two lower-limit closes without instantaneous netting}
\label{sec:conIII}
\begin{theorem}[An internally certified $\mathsf C_1$ witness]\label{thm:two}
Take ten assets with $p=0$, $F=\eta=1$, $K=2$, and sessions
$[0,20]$ and $[20,40]$. There is an absolutely continuous
$\mathsf C_1$ execution such that five targets close at $3/4$
and then at $9/16$, all assets obey sessionwise rolling
$\pm25\%$ bands, all dollar ledgers close at $40$, initial
share inventory $2$ suffices for every asset, and
\begin{equation}\label{eq:certresult}
 M<11.261627<14,\quad \sup_{i,t}r_i^{\rm gross}(t)<0.133,
 \quad G(40)<0.84,
\end{equation}
where $G=(\prod_i x_i)^{1/10}$ is the equal-weight geometric
price index. Hence $\Cris_{1/2}(40)>-\log(.84)>0$ and
$1-G(40)>16\%$.
The parameter and path inequalities are \emph{internally}
verified by outward-rounded interval arithmetic, subject to the
implementation described in Appendix~\ref{app:interval};
external certification is not claimed.
\end{theorem}
\begin{proof}
Partition the ten assets into five disjoint target/compensator
pairs $(x_j,y_j)$. For $K=2$, a pure buyer and pure seller satisfy
\begin{equation}\label{eq:primitives}
 P_+'(s)=\frac{s}{1-2s(s-1)},\qquad
 P_-'(s)=\frac{s}{1+2s(s-1)}.
\end{equation}
Both are strictly increasing on the price range used below.
Buying $z$ dollars sends a price $x$ to the unique $x'$ satisfying
$P_+(x')-P_+(x)=z$; selling $z$ dollars sends $y$ to $y'$ with
$P_-(y)-P_-(y')=z$. These are the \emph{complete} activity ODE
flows, not pure-impact approximations.

In each pair use nine excursions: buy the target while selling the
compensator to the common floor $c$, then (on the first eight
excursions) sell the target to the prescribed rational lows
\begin{equation}\label{eq:lows}
 (1.085409,1.087707,1.085158,1.099042,
 1.106127,1.122661,1.14,1.15).
\end{equation}
After the ninth purchase, sell each target to $3/4$ at the end of
session 1; in session 2 sell it further to $9/16$ while the
corresponding compensator buys. Every amount is determined
uniquely by \eqref{eq:primitives}. Let $f(c)$ be the resulting
signed target-dollar ledger. Directed interval evaluation gives
\begin{equation}\label{eq:root}
 f(0.76065368494358)>0>f(0.76065368494359).
\end{equation}
Because all inverse branches are continuous and strictly monotone,
there is an exact $c_*$ in this bracket with $f(c_*)=0$.
The paired compensator ledger is then also zero.

In the remainder of session 2 sequentially execute the five
compensator-only transfers
$y_1\to y_2\to y_3\to y_4\to y_5\to y_1$, each of
mass exactly $1/25=0.04$. This adds zero net dollars to every
compensator and avoids opposite-direction flows at any instant.
The certificate checks all 19 pair-segment endpoints and all five
cycle transfers; monotonicity between endpoints extends the
rolling-band bounds to the entire continuous path.

Let $A_1,A_2$ be the dollar amounts per pair in sessions 1 and 2.
Five disjoint pairs act in parallel at constant per-edge dollar rate
$A_1/20$ in the first session. In the second, pair transfers followed
by the five cycle edges are sequentially assigned the common
active-edge rate $(A_2+5/25)/20$. Each session lasts exactly 20.
The interval check yields $A_1/20/(3/4)<0.133$ and
$(A_2+1/5)/20/(9/16)<0.133$.
The conservative target inventory is $A/(2\cdot9/16)<2$
where $A=A_1+A_2$; compensator inventory is bounded by
$A/(2c_*)+(1/25)/1.2<2$. Finally every compensator's
terminal price is strictly below $1.23$, so
$G^2<(9/16)\cdot1.23<0.84^2$.
The interval proof details and reproducibility boundary appear in
Appendix~\ref{app:interval}.
\end{proof}

\section{Interpretation and limitations}\label{sec:discussion}
The activity clock breaks the trade-time conservation of
$\sum_i\Phi(S_i)$. Its rate is instead
$-\kappa\sum_iR(S_i)\dd A_i$; the sign depends on the placement
of gross activity across price regions. The local triangular protocol
is already available in $\mathsf C_1$, but the large existence
constructions I--II exploit the broader $\mathsf C_0$ by assigning
common activity even to nodes with small net dollar rates. Such
churning is real executed dollar volume under the stipulated control
rules, not a claim about rational investors' equilibrium behavior.
The $\mathsf C_1$ two-session construction excludes instantaneous
netting but uses many temporal direction reversals; prohibiting
round trips altogether is a different admissibility restriction.

For $p<1$ the scalar floor boundary $K_c(p)$ rules out
unbounded severity at sufficiently large $K$, while Theorem
\ref{thm:nine} establishes a nonempty small-$K$ existence region.
The intervening parameter range and sharp resource-efficient bounds
remain open. Price-floor nonexistence alone is not a sufficiency
theorem. The two-session witness establishes attainability for
\emph{five} targets, not optimality among six, seven or more.
In particular the former seven-target obstruction of the
physical-time model is not transferred here.

\section{Conclusion}
In an idealized fixed-benchmark market where every dollar transfer
is paired and every asset's terminal dollar ledger closes, gross
trade activity can induce both price impact and resilience.
Three-asset trading cycles yield local all-market losses, while
scalar floors, aggregate barriers and necessary execution bounds
restrict their scale. Exact analytic continuation establishes
all-ten halving and a nine-target uniform-resource severity family;
an independent internally certified ten-asset control achieves
five repeated lower-limit closes with finite rate and strict
sessionwise bands. These conclusions concern deterministic
admissibility, not real-world crash probability or strategic
profitability.

\appendix
\section{Continuous-direction construction and exact existence}\label{app:smoothing}
This appendix spells out the smoothing argument used in both
Theorems~\ref{thm:tenhalf} and \ref{thm:nine}; no old instantaneous
pulse is treated as a legal pulse of the new activity model.

\paragraph{Radial identities and bounded-rate interpolation.}
For $p=0$, $F=\eta=1$, choose a nonnegative unit direction
$v(t)$, $\sum_i v_i(t)^2=1$, and set $x_i=R(t)v_i(t)$.
Define $R$ by
\begin{equation}\label{eq:radial}
 \dot R=-\rho\bigl(R-\sum_i v_i\bigr),\quad R(0)=\sqrt N.
\end{equation}
Since $1\le\sum_i v_i\le\sqrt N$, comparison gives
$1\le R\le\sqrt N$. Put $Y_i=v_i^2$.
Differentiating $x_i^2=R^2Y_i$ and using \eqref{eq:radial}
shows that $u_i=x_i[\dot x_i+\rho(x_i-1)]$ obeys
\begin{equation}\label{eq:rateformula}
 u_i=R\dot RY_i+\tfrac12R^2\dot Y_i
      +\rho\bigl(R^2Y_i-R\sqrt{Y_i}\bigr),\qquad \sum_i u_i=0.
\end{equation}
A transition between any two unit directions is realized by
$Y_i(t)=(1-\lambda(t))Y_i^-+\lambda(t)Y_i^+$,
where $\lambda$ increases linearly from zero to one over
width $h>0$; the sum of the squared components remains one.
Each transition is strictly positive for positive endpoints.
Since $|\dot Y_i|\le1/h$, \eqref{eq:rateformula} gives
\begin{equation}\label{eq:a0}
 |u_i|\le \frac{N}{2h}+\rho(2N+\sqrt N)<:a_0,
\quad K=\rho/a_0,\quad M=Na_0T/2.
\end{equation}
For definiteness take $a_0=N/(2h)+\rho(2N+\sqrt N)+1$.
All parameters used to choose $a_0$ are independent of the eventual
terminal target severity $q$.

\paragraph{Convergence of the equation map and its derivatives.}
Take distinct switch times $0<t_1<\cdots<t_{m-1}<T$ with
positive minimum separation. Place an interpolation window
$[t_j-h/2,t_j+h/2]$ around each interior switch and one-sided
windows $[0,h]$ and $[T-h,T]$ at the endpoints. Outside these
windows the direction is one of the fixed stage vectors.
The finite-width direction $v_h$ converges in $L^1$ to the
step direction $v_0$ uniformly as the cut vector and other parameters
range over compact subsets with positive stage gaps. From
\eqref{eq:radial}, $R_h\to R_0$ uniformly. The ledger integrands
$x_{i,h}(x_{i,h}-1)$ also converge in $L^1$.

For derivatives with respect to an interior cut, note that
all stage vectors are permutations of a \emph{single} vector
and have the same coordinate sum $H$. Thus $\sum_i v_0(t)=H$
is constant in cut locations and $\partial_{t_j}R_0=0$.
Every finite-width transition is a translated window with
bounded $L^1$ difference from the sharp step; differentiating its
integrated contribution with respect to its location gives the
jump of the adjoining stage integrands in the zero-width limit.
The remaining effect of the radial equation is $o(1)$:
the signed integral of the change in $\sum_i v_h$ across a
translated window is $O(h)$, since its two flanking stage sums
are the same. Variation of constants in \eqref{eq:radial}
then gives $\partial_{t_j}R_h=O(h)$ away from the $O(h)$
transition window, and its contribution to the ledger derivative
vanishes. On the switch window, the price integrands are uniformly
bounded, so dominated convergence applies after rescaling its
local time coordinate. The endpoint transitions do not move with
interior cuts. Consequently the terminal price and ledger-difference
maps converge in $C^1$ with respect to the cuts. The claim remains
uniform as $q\downarrow0$ in Construction II: the final squared
direction has finite nonnegative components, its price integrands
are bounded continuous functions of $\sqrt{Y_i}$, and its window is
independent of the interior cuts. Only $q>0$ is used for admissible
finished paths.

\paragraph{Invertibility of the cut Jacobian.}
Write the positive high-minus-low weight as $W(t)$ and let
$a_j=\int_{t_{j-1}}^{t_j}W(t)\dd t$ for $j=1,\ldots,m$.
The $m-1$ closure-equations are $a_j-a_m=0$ for $j<m$.
The map $(t_1,\ldots,t_{m-1})\mapsto(a_1,\ldots,a_{m-1})$
is lower triangular with diagonal $W(t_j)>0$.
The map $(a_1,\ldots,a_{m-1})\mapsto(a_j-a_m)_{j<m}$
has derivative $I+\one\one^\top$, determinant $m$.
Thus the exact cut Jacobian determinant is
\begin{equation}\label{eq:cutdet}
 \det\frac{\partial(a_1-a_m,\ldots,a_{m-1}-a_m)}
 {\partial(t_1,\ldots,t_{m-1})}
 =m\prod_{j=1}^{m-1}W(t_j)>0.
\end{equation}
The $C^1$ stability of nondegenerate zeros applies to
Construction I jointly with its monotone terminal-radius equation.
For Construction II it applies uniformly for
$b\in[b_-,\beta]$, because this interval is compact, the cut
weights have a uniform positive lower bound, and the equal-weight
partitions have uniformly positive gaps. Fix one sufficiently small
$h>0$ so that all these conclusions hold simultaneously.
The endpoint sign changes in the compensator ledger then persist
by uniform continuity in $(q,b,h)$, and the intermediate value
theorem determines a compensator-closing $b=b(q)$.

\paragraph{Uniform inventory for arbitrarily small final prices.}
In Construction II, all entries of the stage vectors and the
initial direction are bounded below by
$d=\min(c,b_-,1/\sqrt{10})>0$; the terminal direction can have
nine entries equal to $q\downarrow0$. Before the last window,
$x_i\ge d$. Within the last width-$h$ window, with normalized
transition position $\lambda$, $Y_i\ge(1-\lambda)d^2$ and
$x_i\ge d\sqrt{1-\lambda}$. Thus
\begin{equation}\label{eq:uniformstock}
 \int_0^T\frac{\dd t}{x_i(t)}
 \le\frac{T-h}{d}+\frac{2h}{d}=\frac{T+h}{d},\quad
 H_i^-\le a_0\frac{T+h}{d}.
\end{equation}
Gross-dollar seller rates $s_i=(a_0-u_i)/2$ are at most $a_0$,
and the estimate includes all simultaneous offsetting flow.
Both $M_0=5a_0T$ and the inventory bound are fixed independent
of $q$. All $q>0$ paths are strictly positive-price and
absolutely continuous; no infinite-speed control is used.

\section{Independent interval checker for Construction III}\label{app:interval}
The executable is
\texttt{verification/activity\_driven/certify\_five\_targets.py}.
It does not use V16's physical-time ODE or its 80,000-step
certificate. Its price primitive is integrated from
\eqref{eq:primitives}; explicit antiderivatives are
\begin{align*}
 P_+(s)&=-\frac14\log\bigl(\tfrac32-2(s-\tfrac12)^2\bigr)
 +\frac1{2\sqrt3}\operatorname{artanh}
   \!\left(\frac{2(s-\tfrac12)}{\sqrt3}\right),\\
 P_-(s)&=\frac14\log\bigl(2(s-\tfrac12)^2+\tfrac12\bigr)
 +\frac12\arctan(2s-1).
\end{align*}
The program computes these with outward-rounded interval
operations; the arctangent uses a finite alternating Taylor series
with a separately included geometric remainder after the
argument reduction $\arctan(t)=\pi/4+\arctan((t-1)/(t+1))$.
Each inverse is enclosed by rational endpoints whose primitive
interval evaluations lie strictly below and above the desired
primitive interval. The root is isolated by strict interval signs
at \eqref{eq:root}; all path prices for the resulting exact root
are then propagated as intervals. Since buying and selling have
strictly monotone price flows on the certified ranges, checking
all event endpoints proves the sessionwise bounds between events.
The checker then verifies interval upper bounds on $M$, the
per-asset dollar sales and conservative share inventories, the
specified physical-time rate schedule, and the elementary exact
inequality $(9/16)1.23<0.84^2$. It prints a JSON evidence record.
The JSON records the certified rate schedule separately from a
historical floating-point trial rate of $0.28$; these are not
interchangeable. The checker's scope is an \emph{internally
implemented} interval proof, not a formally verified implementation
or external independent peer review.

\bibliographystyle{unsrtnat}
\enlargethispage{2\baselineskip}
{\small\bibliography{references}}
\end{document}
